\documentclass[11pt,a4paper]{article}

% Compila en Overleaf con pdfLaTeX (configuración por defecto).
\usepackage[utf8]{inputenc}
\usepackage[T1]{fontenc}
\usepackage[spanish,es-noshorthands]{babel}
\usepackage{lmodern}
\usepackage[margin=2.2cm]{geometry}
\usepackage{booktabs,longtable,array,tabularx}
\usepackage[table]{xcolor}
\usepackage{tcolorbox}
\tcbuselibrary{breakable,skins}
\usepackage{enumitem}
\usepackage{hyperref}
\usepackage{fancyhdr}

\definecolor{primario}{HTML}{006B5E}
\definecolor{suave}{HTML}{E6F2EF}
\definecolor{alerta}{HTML}{B3261E}
\definecolor{gris}{HTML}{5F6B68}

\hypersetup{colorlinks=true,linkcolor=primario,urlcolor=primario}
\setlist{itemsep=2pt,topsep=4pt}
\renewcommand{\arraystretch}{1.25}

\pagestyle{fancy}
\fancyhf{}
\lhead{\small Motor de reglas de cartera · Financiera}
\rhead{\small Priorización de historias de usuario}
\cfoot{\small\thepage}

% ── Gherkin ──────────────────────────────────────────────────────────────────
\newtcolorbox{gherkin}{breakable,enhanced,colback=suave,colframe=primario,
  boxrule=0.4pt,arc=2pt,left=6pt,right=6pt,top=4pt,bottom=4pt,fontupper=\small\ttfamily}
\newcommand{\kw}[1]{\textbf{\color{primario}#1}}
\newcommand{\nv}[1]{\par\hspace*{#1em}}

% ── Historia de usuario ──────────────────────────────────────────────────────
\newcommand{\hu}[5]{%
  \subsection*{#1 · #2}
  \noindent\textbf{Prioridad:} #3 \quad\textbf{Épica:} #4\par\smallskip
  \noindent #5\par\smallskip}
\newcommand{\como}[3]{\textbf{Como} #1, \textbf{quiero} #2, \textbf{para} #3.}
\newcommand{\paraht}[2]{\textbf{Para} #1, \textbf{se necesita} #2.}
\newcommand{\nota}[1]{\par\smallskip\noindent{\small\color{gris}\textit{Nota: #1}}}

\newcommand{\must}{\colorbox{alerta}{\color{white}\textbf{\small Must}}}
\newcommand{\should}{\colorbox{primario}{\color{white}\textbf{\small Should}}}
\newcommand{\could}{\colorbox{gris}{\color{white}\textbf{\small Could}}}
\newcommand{\wont}{\colorbox{gray!40}{\textbf{\small Won't (ahora)}}}

\title{\color{primario}\textbf{Motor de reglas de cartera de motos}\\[4pt]
\large Contrato de capacidad y priorización de historias de usuario}
\author{Proyecto Financiera · Azure DevOps \texttt{motosdelcaribe/Financiera}}
\date{8 de octubre de 2026}

\begin{document}
\maketitle
\thispagestyle{fancy}

\begin{tcolorbox}[colback=white,colframe=gris,boxrule=0.3pt,arc=2pt]
\small
\textbf{Cómo se hizo.} Se aplicó la skill \texttt{product-capability} del repositorio ECC
(\texttt{affaan-m/ECC}): primero se fija qué debe ser cierto antes de programar (capacidad,
restricciones, contrato, lo que no se hace, preguntas abiertas) y de ahí sale el orden de trabajo.
Las reglas vienen del Excel \textit{Casos Operaciones V2 mejorado.xlsx} (Diego González);
las cifras de los escenarios se verificaron contra ese Excel.
Arquitectura decidida: \textbf{monolito modular}. El backlog separa historias de usuario (HU)
e historias técnicas (HT).
\end{tcolorbox}

\tableofcontents
\newpage

% ═════════════════════════════════════════════════════════════════════════════
\section{Capacidad}

Un \textbf{motor de reglas de cartera} para créditos de moto con cuotas semanales que, ante cada
abono registrado en el CRM, identifica solo el caso que aplica (pago normal, pago inferior, pago
tardío, abono extra, prepago, cambio de tasa, dación o retoma), lo aplica con las reglas del Excel
y genera el asiento contable en partida doble. El front de los asesores solo origina el crédito y
muestra el plan de la operación normal (caso~1). Lo demás ocurre entre el CRM y el motor, sin
pantallas.

\begin{itemize}
  \item \textbf{Actor principal:} el CRM (sistema) y, detrás, el área de cartera y contabilidad.
  \item \textbf{Resultado después de salir:} ningún abono se aplica ni se contabiliza a mano.
  \item \textbf{Señal de éxito:} las 33 pruebas de la hoja \textit{Validaciones} pasan con
        tolerancia de 1~COP y los asientos cuadran el 100~\% de las veces.
\end{itemize}

% ═════════════════════════════════════════════════════════════════════════════
\section{Restricciones}

\subsection*{Reglas de negocio fijas (del Excel)}
\begin{itemize}
  \item Cuotas semanales: 52 semanas/año, $\lfloor \text{meses}\times 52/12 \rfloor$ cuotas
        (20 meses $\Rightarrow$ 86 cuotas).
  \item Tasa semanal $= (1+\text{EA})^{1/52}-1$; amortización francesa (PMT).
  \item Monto financiado $=$ moto $+$ alistamiento $+$ tecnología $+$ gastos comerciales.
  \item Servicios por cuota: FCC (3~\% mensual $+$ IVA), asistencia (40.000/mes $+$ IVA),
        seguro de vida (0,25~\% mensual, sin IVA).
  \item Prelación (según fórmulas): cobranza $\rightarrow$ mora $\rightarrow$ interés corriente
        $\rightarrow$ servicios $\rightarrow$ capital. \textit{Pendiente de confirmar.}
\end{itemize}

\subsection*{Invariantes (nunca se pueden romper)}
\begin{enumerate}
  \item El saldo de capital nunca es negativo.
  \item Todo asiento cuadra: débitos $=$ créditos (tolerancia 1~COP).
  \item La suma de la distribución de un abono es igual al valor del abono.
  \item Un mismo abono (mismo identificador) nunca se aplica dos veces.
  \item Un crédito cancelado o cerrado por evento no recibe abonos.
  \item Ningún movimiento se borra: los errores se corrigen con reversos.
  \item Cada movimiento guarda la versión de reglas y de parámetros con que se calculó.
\end{enumerate}

\subsection*{Restricciones regulatorias y técnicas}
\begin{itemize}
  \item La tasa no supera la usura de la modalidad; si la supera: usura $-$ 1 punto.
  \item Datos personales según la Ley 1581 de 2012.
  \item Cálculo con decimales exactos (nunca punto flotante).
  \item Arquitectura: monolito modular (un servicio, una base de datos, módulos con fronteras).
\end{itemize}

% ═════════════════════════════════════════════════════════════════════════════
\section{Contrato de implementación}

\subsection*{Actores y superficies}
\begin{tabularx}{\linewidth}{>{\bfseries}l X}
\toprule
Actor & Superficie \\
\midrule
Asesor del concesionario & Front React: cotizar, originar, plan de la operación normal. \\
CRM & API: registrar abono, consultar crédito; recibe avisos de cambio de estado. \\
Contabilidad & Exportación de asientos y reporte de terceros. \\
Procesos programados & Cierre diario, revisión de usura, liquidación mensual. \\
\bottomrule
\end{tabularx}

\subsection*{Estados del crédito y transiciones}
\begin{center}
\begin{tabular}{l c l l}
\toprule
Desde & & Hacia & Evento \\
\midrule
Originado & $\rightarrow$ & Vigente al día & Desembolso \\
Vigente al día & $\rightarrow$ & En mora & Pago inferior o cuota vencida sin pago \\
En mora & $\rightarrow$ & Vigente al día & Pagos que cubren todo lo vencido \\
Vigente al día & $\rightarrow$ & Cancelado & Prepago total o última cuota pagada \\
En mora & $\rightarrow$ & Cancelado & Pago de todo lo adeudado \\
En mora & $\rightarrow$ & Cerrado por evento & Dación en pago o retoma \\
\bottomrule
\end{tabular}
\end{center}

\subsection*{Módulos del monolito}
\begin{tabularx}{\linewidth}{>{\ttfamily}l X}
\toprule
\normalfont\bfseries Módulo & \bfseries Responsabilidad \\
\midrule
nucleo-calculo & Reglas puras de los casos 1 a 9; sin base de datos ni HTTP. \\
creditos & Originación, historial de movimientos, estados, versiones. \\
abonos & Clasificador del caso, prelación, idempotencia, reversos. \\
contabilidad & Plan de cuentas, asientos, cuadre, exportación. \\
integracion-crm & API, avisos con reintentos (outbox), conciliación. \\
procesos & Cierre diario, usura, liquidación a terceros. \\
api & Autenticación, validación de entrada, controladores. \\
\bottomrule
\end{tabularx}

% ═════════════════════════════════════════════════════════════════════════════
\section{Lo que el motor no hace}
\begin{itemize}
  \item No aprueba créditos (estudio, centrales, visita domiciliaria): eso es del CRM y del área de riesgo.
  \item No recauda dinero ni se conecta con bancos o PSE: recibe abonos ya recaudados.
  \item No se comunica con el cliente (mensajes, WhatsApp, llamadas): emite avisos al CRM.
  \item No es el sistema contable oficial: genera los asientos y los exporta.
  \item No tiene pantallas para los casos 2 a 9.
\end{itemize}

% ═════════════════════════════════════════════════════════════════════════════
\section{Preguntas abiertas que bloquean}

\begin{longtable}{>{\bfseries}p{0.6cm} p{8.6cm} p{5.6cm}}
\toprule
\# & Pregunta & Bloquea \\
\midrule
\endhead
P1 & ¿Prelación oficial? ¿Servicios antes o después del capital? & Casos 8 y 9, contabilidad \\
P2 & ¿Se calcula con centavos o se redondea a pesos? ¿Dónde se absorbe el residuo? & Todo el núcleo \\
P3 & ¿Fórmula única de mora (semanal o por días)? & Casos 8 y 9 \\
P4 & ¿Hay cuota inicial? ¿LTV máximo? & Originación y front \\
P5 & ¿La EA de 87,32~\% cumple la usura de la modalidad? & Salida a producción \\
P6 & ¿FCC y seguro sobre el monto inicial o sobre el saldo? & Casos 5, 6 y 7 \\
P7 & ¿Mora en dación y retoma? & Casos 3 y 4 \\
P8 & ¿Qué CRM es y cómo se integra (API, webhooks, n8n)? & Integración \\
P9 & ¿Cuál es el sistema contable oficial y su formato? & Exportación contable \\
P10 & ¿Abono tardío e incompleto a la vez (caso 10)? & Clasificador \\
\bottomrule
\end{longtable}

% ═════════════════════════════════════════════════════════════════════════════
\section{Historias de usuario e historias técnicas}

El backlog separa dos tipos de trabajo:

\begin{tabularx}{\linewidth}{>{\bfseries}l X X}
\toprule
 & Historia de usuario (HU) & Historia técnica (HT) \\
\midrule
Qué es & Valor que un usuario o el negocio puede ver y validar. & Trabajo que el sistema necesita para que las HU funcionen, sean confiables o se puedan probar. \\
Formato & \textbf{Como} \ldots, \textbf{quiero} \ldots, \textbf{para} \ldots & \textbf{Para} \ldots, \textbf{se necesita} \ldots \\
Se demuestra & En la revisión del sprint, a negocio. & Con evidencia técnica: pruebas, pipeline, ADR. \\
En Azure DevOps & User Story / Product Backlog Item. & Mismo tipo, con prefijo \texttt{[HT]} y etiqueta \texttt{historia-tecnica}. \\
\bottomrule
\end{tabularx}

\subsection*{Historias que pasan a ser técnicas}
\begin{tabularx}{\linewidth}{>{\bfseries}l >{\bfseries}l X}
\toprule
Antes & Ahora & Historia \\
\midrule
HU-06 & HT-01 & ADR de monolito modular y esqueleto de los módulos \\
HU-18 & HT-02 & Proteger los invariantes del motor \\
HU-20 & HT-03 & Historial inmutable de movimientos y recálculo \\
HU-21 & HT-04 & Versionar parámetros con fecha de vigencia \\
HU-24 & HT-05 & Evitar aplicar dos veces el mismo abono (idempotencia) \\
HU-31 & HT-06 & Rechazar asientos descuadrados \\
HU-37 & HT-07 & Avisos al CRM sin pérdida (outbox) \\
\bottomrule
\end{tabularx}

\subsection*{Historias técnicas nuevas}
\begin{tabularx}{\linewidth}{>{\bfseries}l X l}
\toprule
ID & Historia & Sprint \\
\midrule
HT-08 & Repositorio y estructura del proyecto & 1 \\
HT-09 & Pipeline de integración continua con pruebas y prueba de arquitectura & 1 \\
HT-10 & Cronogramas del Excel como datos de prueba & 1 \\
HT-11 & Ambiente de pruebas desplegado & 3 \\
\bottomrule
\end{tabularx}

% ═════════════════════════════════════════════════════════════════════════════
\section{Plan de 3 sprints de una semana}

\subsection*{Criterio}
Se ordena por: (1) lo que \textbf{bloquea} a otras historias va primero; (2) lo que
\textbf{reduce el riesgo} de calcular mal el dinero; (3) el \textbf{valor} para el negocio;
(4) el tamaño. La prioridad usa MoSCoW: \must{} imprescindible en el sprint,
\should{} entra si alcanza la capacidad, \could{} deseable, \wont{} fuera de estos 3 sprints.
Las historias técnicas se priorizan igual que las de usuario: van cuando otra historia las necesita.

\begin{tcolorbox}[colback=white,colframe=alerta,boxrule=0.4pt,arc=2pt]
\small\textbf{Alcance realista.} Los 3 sprints suman \textbf{98 puntos}. Con el supuesto de
\textbf{25 puntos por semana} caben unos 75: sobran cerca de 23. Si la velocidad real del Sprint~1
confirma ese ritmo, se abre un Sprint~4 y se mueven, en este orden: HU-39, HT-04, HT-11, HU-36,
HU-38 y la contabilidad base (HU-29, HU-30, HT-06).
Estos sprints entregan el \textbf{núcleo del motor}: caso 1, abonos al día y extraordinarios
(casos 5, 6 y 7), mora y cobranza (casos 8 y 9), con contabilidad y avisos al CRM. El cambio de tasa,
la dación, la retoma, la usura automática y el front semanal quedan para después.
\end{tcolorbox}

\begin{longtable}{>{\bfseries}p{1.4cm} p{6.4cm} p{1.2cm} p{1.2cm} p{1.2cm} p{2.6cm}}
\toprule
Sprint & Objetivo (lo que se demuestra el viernes) & HU & HT & Puntos & Casos \\
\midrule
\endhead
Sprint 1 & Reglas cerradas con negocio y núcleo de cálculo que reproduce el caso 1 al centavo & 13 & 5 & 29 & Caso 1 \\
Sprint 2 & El motor recibe abonos al día y extraordinarios por API, los aplica y los contabiliza & 10 & 4 & 35 & 1, 5, 6 y 7 \\
Sprint 3 & Pagos incompletos, mora y cobranza, con avisos al CRM sin pérdida & 9 & 2 & 34 & 8 y 9 \\
Después & Cambio de tasa, usura, dación, retoma, front semanal, terceros & 10 & — & — & 2, 3 y 4 \\
\bottomrule
\end{longtable}

% ── Sprint 1 ─────────────────────────────────────────────────────────────────
\subsection{Sprint 1 · Semana 1 (priorizado)}

\noindent\textbf{Objetivo:} que negocio cierre las reglas que bloquean, que el proyecto quede
montado con su pipeline y que el núcleo de cálculo reproduzca el caso~1 del Excel con tolerancia de 1~COP.

\begin{longtable}{>{\bfseries}p{0.6cm} p{1.2cm} p{7.4cm} p{1.5cm} p{1.1cm} p{2.2cm}}
\toprule
\# & ID & Historia & Prioridad & Puntos & Depende de \\
\midrule
\endhead
1 & HU-02 & Resolver la prelación de abonos (P1) & \must & 1 & — \\
2 & HU-03 & Definir el redondeo del motor (P2) & \must & 1 & — \\
3 & HU-04 & Definir la fórmula de mora (P3) & \must & 1 & — \\
4 & HU-05 & Decidir cuota inicial y monto financiado (P4) & \must & 1 & — \\
5 & HU-09 & Confirmar el IVA de los servicios & \must & 1 & — \\
6 & HT-08 & Repositorio y estructura del proyecto & \must & 2 & — \\
7 & HT-01 & ADR de monolito modular y esqueleto de los módulos & \must & 3 & HT-08 \\
8 & HT-09 & Pipeline de integración continua & \must & 2 & HT-08 \\
9 & HT-10 & Cronogramas del Excel como datos de prueba & \must & 1 & — \\
10 & HU-01 & Catálogo de reglas del Excel (caso 1 primero) & \must & 2 & HU-02 a 05 \\
11 & HU-11 & Convertir la EA a tasas semanal, diaria y mensual & \must & 1 & HU-03 \\
12 & HU-12 & Calcular el número de cuotas semanales & \must & 1 & — \\
13 & HU-13 & Calcular el monto financiado por componentes & \must & 1 & HU-05 \\
14 & HU-14 & Calcular la cuota financiera semanal fija & \must & 2 & HU-11, 12 \\
15 & HU-15 & Generar el cronograma semanal de amortización & \must & 3 & HU-14 \\
16 & HU-16 & Calcular los servicios complementarios por cuota & \must & 2 & HU-09 \\
17 & HU-17 & Validaciones del caso 1 (filas 1 a 5 de la hoja Validaciones) & \must & 2 & HU-15, HT-10 \\
18 & HT-02 & Invariantes del cálculo (saldo nunca negativo, cierre en cero) & \should & 2 & HU-15 \\
\midrule
 & & \textbf{Total} (\must{} 27 · \should{} 2) — 13 HU y 5 HT & & \textbf{29} & \\
\bottomrule
\end{longtable}

\noindent{\small\color{alerta}\textit{Pasa de 25 puntos: si no alcanza, sale primero HT-02 (\should) y
luego la validación del FCC del mes~1 dentro de HU-17.}}

\subsubsection*{Plan de la semana}
\begin{tabularx}{\linewidth}{>{\bfseries}l X}
\toprule
Día & Trabajo \\
\midrule
Lunes & Sesión de refinamiento de 2 horas con Diego González y Heider Galvis: cerrar P1 a P4 y el IVA (HU-02 a 05 y HU-09). En paralelo: repositorio (HT-08), ADR y esqueleto del monolito (HT-01). \\
Martes & Pipeline de integración continua (HT-09) y cronogramas del Excel como datos de prueba (HT-10). Catálogo de reglas (HU-01). Tasas, número de cuotas y monto financiado (HU-11 a 13), con pruebas primero (TDD). \\
Miércoles & Cuota financiera (HU-14) y cronograma semanal (HU-15). \\
Jueves & Servicios complementarios (HU-16) e invariantes del cálculo (HT-02). \\
Viernes & Validaciones del caso 1 contra el Excel (HU-17). Revisión del sprint con negocio: se muestra el cronograma de 86 cuotas comparado con el Excel. Retrospectiva. \\
\bottomrule
\end{tabularx}

\subsubsection*{Definición de terminado (Sprint 1)}
\begin{itemize}
  \item Las decisiones P1 a P4 y el IVA quedan escritas y aprobadas por negocio.
  \item Cada regla tiene prueba automática escrita antes del código.
  \item El pipeline corre en cada pull request y está en verde.
  \item Las 5 validaciones del caso~1 pasan con tolerancia de 1~COP: monto desembolsado,
        tasa semanal, cuota financiera, saldo final y pago FCC del mes 1.
  \item El módulo \texttt{nucleo-calculo} no depende de ningún otro módulo (prueba de arquitectura en verde).
  \item Código revisado en pull request y unido a la rama principal.
\end{itemize}

\subsubsection*{Riesgos del Sprint 1}
\begin{itemize}
  \item \textbf{Negocio no cierra P1 a P3 el lunes:} el cálculo del caso~1 no depende de ellas,
        así que se sigue con el núcleo; lo que sí bloquean es el Sprint 3.
  \item \textbf{P2 (redondeo) cambia la convención:} el núcleo se programa con decimales exactos
        y el redondeo en un solo punto, para cambiarlo sin tocar las fórmulas.
  \item \textbf{No está decidido el lenguaje:} HT-08 lo fija el lunes; sin eso no arranca nada más.
\end{itemize}

% ── Sprint 2 y 3 ─────────────────────────────────────────────────────────────
\subsection{Sprint 2 · Semana 2}
\noindent\textbf{Objetivo:} el motor recibe abonos al día y extraordinarios por API, los aplica y los contabiliza.

\begin{longtable}{>{\bfseries}p{1.2cm} p{8.6cm} p{1.5cm} p{1.1cm} p{2.2cm}}
\toprule
ID & Historia & Prioridad & Puntos & Depende de \\
\midrule
\endhead
HU-07 & Contrato de datos con el CRM & \must & 2 & — \\
HU-19 & Máquina de estados del crédito & \must & 3 & — \\
HT-03 & Historial inmutable de movimientos y recálculo & \must & 5 & HT-01 \\
HT-04 & Versionar parámetros con fecha de vigencia & \must & 2 & HT-03 \\
HU-23 & API para registrar abonos desde el CRM & \must & 3 & HU-07, HT-03 \\
HT-05 & Evitar aplicar dos veces el mismo abono & \must & 2 & HU-23 \\
HU-25 & Clasificar el abono (normal, extra, prepago) & \must & 3 & HU-19 \\
HU-26 & Prepago total (caso 5) & \must & 2 & HU-25 \\
HU-27 & Abono extra con menor plazo (caso 6) & \must & 2 & HU-25 \\
HU-28 & Abono extra con menor cuota (caso 7) & \must & 2 & HU-25 \\
HU-29 & Configurar el mapeo de conceptos a cuentas & \must & 2 & — \\
HU-30 & Contabilizar la cuota normal y los abonos extra & \must & 3 & HU-29 \\
HT-06 & Rechazar asientos descuadrados & \must & 1 & HU-30 \\
HU-17 & Validaciones de los casos 5, 6 y 7 & \must & 3 & HU-26 a 28 \\
\midrule
 & \textbf{Total} — 10 HU y 4 HT & & \textbf{35} & \\
\bottomrule
\end{longtable}

\subsection{Sprint 3 · Semana 3}
\noindent\textbf{Objetivo:} pagos incompletos, mora y cobranza, con avisos al CRM sin pérdida.

\begin{longtable}{>{\bfseries}p{1.2cm} p{8.6cm} p{1.5cm} p{1.1cm} p{2.2cm}}
\toprule
ID & Historia & Prioridad & Puntos & Depende de \\
\midrule
\endhead
HU-22 & Calcular fechas de vencimiento semanales & \must & 2 & HU-15 \\
HU-32 & Calcular días de atraso & \must & 2 & HU-22 \\
HU-33 & Aplicar un pago inferior con la prelación (caso 8) & \must & 5 & HU-02, 25 \\
HU-34 & Causar interés de mora sobre capital vencido & \must & 3 & HU-04 \\
HU-35 & Gastos de cobranza por días de atraso (caso 9) & \must & 3 & HU-32 \\
HU-36 & Cierre diario de cartera & \must & 3 & HU-34 \\
HT-07 & Avisos al CRM sin pérdida (outbox) & \must & 3 & HU-19 \\
HT-11 & Ambiente de pruebas desplegado & \must & 3 & HT-09 \\
HU-38 & Contabilizar pagos inferiores, mora y cobranza & \must & 5 & HU-33 \\
HU-17 & Validaciones de los casos 8 y 9 & \must & 3 & HU-33 a 35 \\
HU-39 & Reversar un abono aplicado por error & \should & 2 & HT-03 \\
\midrule
 & \textbf{Total} — 9 HU y 2 HT & & \textbf{34} & \\
\bottomrule
\end{longtable}

\subsection{Después de los 3 sprints}
\begin{longtable}{>{\bfseries}p{1.2cm} p{10.2cm} p{1.5cm}}
\toprule
ID & Historia & Prioridad \\
\midrule
\endhead
HU-08 & Validar la tasa EA frente a la usura (debe estar antes de producción) & \should \\
HU-40 & Cambio de tasa (caso 2) & \should \\
HU-41 & Ajuste automático por usura & \should \\
HU-42 & Dación en pago (caso 3) y retoma (caso 4) & \should \\
HU-43 & Cotizador semanal en el front (llama al motor) & \should \\
HU-44 & Reporte mensual de pagos a terceros & \should \\
HU-45 & Conciliación diaria CRM–motor & \should \\
HU-10 & Spike de Fineract (solo para confirmar la decisión) & \could \\
HU-46 & Combinar casos: abono tardío e incompleto (P10) & \could \\
HU-47 & Proyección al horizonte de 120 semanas & \wont \\
\bottomrule
\end{longtable}

% ═════════════════════════════════════════════════════════════════════════════
\section{Sprint 1 en detalle}

\textit{Criterios de aceptación de las 18 historias del Sprint 1 (13 HU y 5 HT), en el orden en que
se trabajan. Las nuevas de los sprints 2 y 3 van en la sección siguiente; el resto conserva los
criterios del backlog cargado en Azure DevOps.}

% ── Decisiones (lunes) ───────────────────────────────────────────────────────
\hu{HU-02}{Resolver la prelación de abonos}{\must}{E1 · Descubrimiento}{%
\como{Product Owner}{confirmar con negocio el orden exacto en que se aplica un abono}{programar una sola prelación sin ambigüedad}}
\begin{gherkin}
\kw{Característica:} Prelación oficial de aplicación de abonos
\nv{1}\kw{Escenario:} Se elimina la contradicción entre texto y fórmulas
\nv{2}\kw{Dado} que la hoja Casos describe "mora, interés corriente, capital"
\nv{2}\kw{Y} que las fórmulas aplican "cobranza, mora, interés, servicios, capital"
\nv{2}\kw{Cuando} negocio define la prelación oficial
\nv{2}\kw{Entonces} queda un único orden que incluye todos los conceptos
\nv{2}\kw{Y} el Excel se actualiza para que texto y fórmulas coincidan
\end{gherkin}

\hu{HU-03}{Definir el redondeo del motor}{\must}{E1 · Descubrimiento}{%
\como{Product Owner}{definir si el motor trabaja con centavos o redondea a pesos}{que las pruebas contra el Excel tengan una tolerancia clara}}
\begin{gherkin}
\kw{Característica:} Precisión de los cálculos
\nv{1}\kw{Escenario:} Se fija la convención de redondeo
\nv{2}\kw{Dado} que el Excel no redondea y acepta 1 COP de tolerancia
\nv{2}\kw{Cuando} negocio y contabilidad definen la convención
\nv{2}\kw{Entonces} queda escrito con cuántos decimales se calcula y se guarda
\nv{2}\kw{Y} cómo se redondea lo que se cobra y lo que se contabiliza
\nv{2}\kw{Y} en qué cuota se absorbe la diferencia de redondeo
\end{gherkin}

\hu{HU-04}{Definir la fórmula de mora}{\must}{E1 · Descubrimiento}{%
\como{Product Owner}{una sola fórmula de interés de mora}{no tener dos cálculos distintos entre el caso 8 y el caso 9}}
\begin{gherkin}
\kw{Característica:} Fórmula única de interés de mora
\nv{1}\kw{Escenario:} Se unifica la mora de los casos 8 y 9
\nv{2}\kw{Dado} que el caso 8 calcula mora = capital vencido x tasa semanal
\nv{2}\kw{Y} que el caso 9 suma capital de la cuota x ((1 + tasa diaria)\^{}días - 1)
\nv{2}\kw{Cuando} negocio define la fórmula oficial
\nv{2}\kw{Entonces} existe una única regla con su base, su tasa y su periodo
\nv{2}\kw{Y} se confirma que la tasa de mora no supera el límite legal
\end{gherkin}

\hu{HU-05}{Decidir cuota inicial y monto financiado}{\must}{E1 · Descubrimiento}{%
\como{Product Owner}{decidir si el crédito exige cuota inicial o financia el 116~\% de la moto como el Excel}{alinear el front con el motor}}
\begin{gherkin}
\kw{Característica:} Composición del monto financiado
\nv{1}\kw{Escenario:} Se define si hay cuota inicial
\nv{2}\kw{Dado} que el Excel financia 5.904.400 sobre una moto de 5.100.000
\nv{2}\kw{Y} que el front actual exige una cuota inicial mínima del 10 \%
\nv{2}\kw{Cuando} negocio toma la decisión
\nv{2}\kw{Entonces} queda definido si existe cuota inicial y su porcentaje mínimo
\nv{2}\kw{Y} si se descuenta de la moto o del total financiado
\nv{2}\kw{Y} el LTV máximo permitido
\end{gherkin}

\hu{HU-09}{Confirmar el IVA de los servicios}{\must}{E1 · Descubrimiento}{%
\como{contador}{confirmar qué servicios llevan IVA del 19~\% y cuáles no}{liquidar bien lo que se cobra y lo que se paga a terceros}}
\begin{gherkin}
\kw{Característica:} IVA de servicios
\nv{1}\kw{Escenario:} Servicios gravados y excluidos
\nv{2}\kw{Dado} que el Excel grava FCC y asistencia y excluye el seguro de vida
\nv{2}\kw{Cuando} contabilidad confirma el tratamiento tributario
\nv{2}\kw{Entonces} cada servicio queda marcado como gravado o excluido con su tarifa
\nv{2}\kw{Y} queda definido quién responde por el IVA de cada servicio
\end{gherkin}

% ── Base técnica (lunes y martes) ────────────────────────────────────────────
\hu{HT-08}{Repositorio y estructura del proyecto}{\must}{E9 · Plataforma \textit{(técnica, nueva)}}{%
\paraht{que el equipo empiece a programar el lunes sin bloquearse}{un repositorio con el lenguaje, la base de datos local, las convenciones y la estructura de carpetas acordadas}}
\begin{gherkin}
\kw{Característica:} Proyecto listo para desarrollar
\nv{1}\kw{Escenario:} Arranque en un equipo nuevo
\nv{2}\kw{Dado} un desarrollador que clona el repositorio
\nv{2}\kw{Cuando} sigue las instrucciones del README
\nv{2}\kw{Entonces} compila el proyecto y corre las pruebas en menos de 15 minutos
\nv{2}\kw{Y} la base de datos local se levanta con un solo comando
\nv{1}\kw{Escenario:} Convenciones escritas
\nv{2}\kw{Dado} el repositorio creado
\nv{2}\kw{Cuando} se revisa la documentación
\nv{2}\kw{Entonces} están el lenguaje y su versión, el estilo de código, las ramas y el flujo de pull request
\end{gherkin}

\hu{HT-01}{ADR de monolito modular y esqueleto de los módulos}{\must}{E9 · Plataforma \textit{(técnica)}}{%
\paraht{construir sobre una base acordada y que las fronteras entre módulos no se rompan}{un registro de decisión (ADR) y el esqueleto de los 7 módulos con sus reglas de dependencia}}
\begin{gherkin}
\kw{Característica:} Arquitectura de monolito modular
\nv{1}\kw{Escenario:} Se documenta la decisión
\nv{2}\kw{Dado} las opciones evaluadas: monolito modular, microservicios y Fineract
\nv{2}\kw{Cuando} el equipo registra el ADR
\nv{2}\kw{Entonces} queda la decisión, sus razones y cuándo separar un módulo
\nv{1}\kw{Escenario:} Las fronteras se verifican solas
\nv{2}\kw{Dado} el módulo nucleo-calculo
\nv{2}\kw{Cuando} corre la prueba de arquitectura
\nv{2}\kw{Entonces} falla si nucleo-calculo depende de otro módulo
\nv{2}\kw{Y} falla si un módulo lee las tablas de otro módulo
\end{gherkin}

\hu{HT-09}{Pipeline de integración continua}{\must}{E9 · Plataforma \textit{(técnica, nueva)}}{%
\paraht{detectar en cada cambio si el motor deja de cuadrar con el Excel}{un pipeline que compile y corra todas las pruebas en cada pull request}}
\begin{gherkin}
\kw{Característica:} Integración continua
\nv{1}\kw{Escenario:} Pull request con pruebas en verde
\nv{2}\kw{Dado} un pull request hacia la rama principal
\nv{2}\kw{Cuando} se ejecuta el pipeline
\nv{2}\kw{Entonces} corren las pruebas unitarias, las del Excel y la de arquitectura
\nv{2}\kw{Y} el resultado queda visible en el pull request
\nv{1}\kw{Escenario:} Una prueba falla
\nv{2}\kw{Dado} un cambio que rompe una validación del Excel
\nv{2}\kw{Cuando} corre el pipeline
\nv{2}\kw{Entonces} el pull request no se puede unir a la rama principal
\end{gherkin}

\hu{HT-10}{Cronogramas del Excel como datos de prueba}{\must}{E8 · Calidad \textit{(técnica, nueva)}}{%
\paraht{comparar el motor contra el Excel sin copiar cifras a mano}{exportar los parámetros y los cronogramas de los 9 casos a archivos que las pruebas puedan leer}}
\begin{gherkin}
\kw{Característica:} Datos de prueba desde el Excel
\nv{1}\kw{Escenario:} Exportación de los casos
\nv{2}\kw{Dado} el Excel "Casos Operaciones V2 mejorado.xlsx"
\nv{2}\kw{Cuando} se ejecuta el exportador
\nv{2}\kw{Entonces} se genera un archivo por caso con todas las filas y columnas del cronograma
\nv{2}\kw{Y} un archivo con las 33 filas de la hoja Validaciones
\nv{1}\kw{Escenario:} Nueva versión del Excel
\nv{2}\kw{Dado} que Diego entrega una versión corregida del Excel
\nv{2}\kw{Cuando} se vuelve a ejecutar el exportador
\nv{2}\kw{Entonces} los archivos se regeneran y las pruebas muestran qué cambió
\end{gherkin}

% ── Núcleo de cálculo (martes a viernes) ─────────────────────────────────────
\hu{HU-01}{Catálogo de reglas del Excel}{\must}{E1 · Descubrimiento}{%
\como{Product Owner}{un catálogo de reglas numeradas extraído de las hojas Inputs y Casos}{que desarrollo, pruebas y negocio hablen de la misma regla}}
\begin{gherkin}
\kw{Característica:} Catálogo de reglas del motor de cartera
\nv{1}\kw{Escenario:} Cada fórmula queda como regla trazable
\nv{2}\kw{Dado} el Excel con sus 12 hojas
\nv{2}\kw{Cuando} se documenta el catálogo
\nv{2}\kw{Entonces} cada fórmula tiene identificador, descripción, fórmula y celda de origen
\nv{2}\kw{Y} indica en qué caso (1 a 9) aplica
\nv{1}\kw{Escenario:} Aprobación por negocio
\nv{2}\kw{Dado} el catálogo documentado
\nv{2}\kw{Cuando} Diego González lo revisa
\nv{2}\kw{Entonces} cada regla queda "aprobada", "ajustar" o "descartar"
\nv{2}\kw{Y} solo se programan reglas aprobadas
\end{gherkin}
\nota{en el Sprint~1 se cataloga primero el caso~1; los demás casos se completan antes de su sprint.}

\hu{HU-11}{Convertir la EA a tasas semanal, diaria y mensual}{\must}{E2 · Originación}{%
\como{analista de crédito}{que el motor derive las tasas periódicas de la EA pactada}{usar la tasa contractual en todos los cálculos}}
\begin{gherkin}
\kw{Característica:} Conversión de tasas
\nv{1}\kw{Escenario:} Tasas equivalentes de una EA de 87,32 \%
\nv{2}\kw{Dado} una tasa efectiva anual de 87,32 \%
\nv{2}\kw{Cuando} se calculan las tasas equivalentes
\nv{2}\kw{Entonces} la tasa semanal es 0,012143296 (tolerancia 0,000001)
\nv{2}\kw{Y} la tasa diaria es 0,001721063
\nv{2}\kw{Y} la tasa mensual de referencia es 0,053696035
\end{gherkin}

\hu{HU-12}{Calcular el número de cuotas semanales}{\must}{E2 · Originación}{%
\como{analista de crédito}{convertir el plazo en meses a cuotas semanales}{generar el cronograma semanal}}
\begin{gherkin}
\kw{Característica:} Número de cuotas semanales
\nv{1}\kw{Esquema del escenario:} Conversión de meses a semanas
\nv{2}\kw{Dado} un año de 52 semanas y 12 meses
\nv{2}\kw{Cuando} el plazo es de <meses> meses
\nv{2}\kw{Entonces} el número de cuotas es <cuotas>
\nv{2}\kw{Ejemplos:}
\nv{3}| meses | cuotas |
\nv{3}| 20 \ \ | 86 \ \ \ |
\nv{3}| 12 \ \ | 52 \ \ \ |
\nv{3}| 18 \ \ | 78 \ \ \ |
\end{gherkin}

\hu{HU-13}{Calcular el monto financiado por componentes}{\must}{E2 · Originación}{%
\como{analista de crédito}{que el monto financiado sume moto, alistamiento, tecnología y gastos comerciales}{colocar el valor final correcto}}
\begin{gherkin}
\kw{Característica:} Monto financiado
\nv{1}\kw{Escenario:} Suma de componentes
\nv{2}\kw{Dado} una moto de 5.100.000
\nv{2}\kw{Y} matrícula y RUNT 171.100, SOAT 373.300 y prenda 120.000
\nv{2}\kw{Y} GPS y plataforma 80.000 y gastos comerciales 60.000
\nv{2}\kw{Cuando} se calcula el monto financiado
\nv{2}\kw{Entonces} el monto financiado es 5.904.400
\nv{2}\kw{Y} el porcentaje de financiación sobre la moto es 115,77 \%
\end{gherkin}

\hu{HU-14}{Calcular la cuota financiera semanal fija}{\must}{E2 · Originación}{%
\como{analista de crédito}{calcular la cuota fija de capital más interés}{informar al cliente cuánto paga cada semana}}
\begin{gherkin}
\kw{Característica:} Cuota financiera
\nv{1}\kw{Escenario:} Cuota del crédito base
\nv{2}\kw{Dado} un monto financiado de 5.904.400
\nv{2}\kw{Y} una tasa semanal de 0,012143296 y 86 cuotas
\nv{2}\kw{Cuando} se calcula la cuota con amortización francesa
\nv{2}\kw{Entonces} la cuota financiera es 111.014,82 (tolerancia 1 COP)
\nv{2}\kw{Y} la cuota total con servicios es 174.048,99
\end{gherkin}

\hu{HU-15}{Generar el cronograma semanal de amortización}{\must}{E2 · Originación}{%
\como{analista de crédito}{el plan cuota a cuota con saldo, interés, capital y servicios}{conocer el comportamiento esperado del crédito}}
\begin{gherkin}
\kw{Característica:} Cronograma de amortización
\nv{1}\kw{Escenario:} Primera cuota
\nv{2}\kw{Dado} el crédito base de 5.904.400 a 86 cuotas
\nv{2}\kw{Cuando} se genera el cronograma
\nv{2}\kw{Entonces} la cuota 1 tiene interés 71.698,88 y capital 39.315,95
\nv{2}\kw{Y} saldo final 5.865.084,05
\nv{1}\kw{Escenario:} Cierre del crédito
\nv{2}\kw{Dado} el mismo cronograma
\nv{2}\kw{Cuando} se llega a la cuota 86
\nv{2}\kw{Entonces} el saldo final es 0 (tolerancia 1 COP)
\nv{2}\kw{Y} los intereses totales son 3.642.874,68
\end{gherkin}

\hu{HU-16}{Calcular los servicios complementarios por cuota}{\must}{E2 · Originación}{%
\como{analista de crédito}{configurar FCC, asistencia y seguro con su base, periodicidad e IVA}{calcular lo que paga el cliente en cada cuota}}
\begin{gherkin}
\kw{Característica:} Servicios complementarios
\nv{1}\kw{Escenario:} Servicios de una cuota semanal
\nv{2}\kw{Dado} un monto financiado de 5.904.400 y 4,333333 semanas por mes
\nv{2}\kw{Cuando} se calculan los servicios por cuota
\nv{2}\kw{Entonces} el FCC con IVA es 48.643,17
\nv{2}\kw{Y} la asistencia con IVA es 10.984,62
\nv{2}\kw{Y} el seguro de vida es 3.406,38
\nv{2}\kw{Y} los servicios por cuota suman 63.034,17
\end{gherkin}

\hu{HU-17}{Automatizar las 33 pruebas de la hoja Validaciones}{\must}{E8 · Calidad}{%
\como{QA}{que cada fila de Validaciones sea una prueba automática}{saber en cada cambio si el motor sigue cuadrando con el Excel}}
\begin{gherkin}
\kw{Característica:} Pruebas de validación del Excel
\nv{1}\kw{Esquema del escenario:} Prueba de la hoja Validaciones
\nv{2}\kw{Dado} el caso <caso> con los parámetros de la hoja Inputs
\nv{2}\kw{Cuando} el motor calcula <variable> en el periodo <periodo>
\nv{2}\kw{Entonces} coincide con el Excel dentro de <tolerancia>
\nv{2}\kw{Ejemplos:}
\nv{3}| Caso 1 | Cuota financiera \ \ \ | 1 \ | 1 COP \ |
\nv{3}| Caso 2 | Nueva cuota \ \ \ \ \ \ \ \ | 44 | 1 COP \ |
\nv{3}| Caso 6 | Último periodo pago | 61 | exacto |
\nv{3}| Caso 8 | Interés mora \ \ \ \ \ \ \ | 44 | 1 COP \ |
\end{gherkin}
\nota{en el Sprint~1 solo se automatizan las 5 validaciones del caso~1; las de los casos 5, 6 y 7
van en el Sprint~2 y las de los casos 8 y 9 en el Sprint~3.}

\hu{HT-02}{Proteger los invariantes del motor}{\should}{E2 · Originación \textit{(técnica)}}{%
\paraht{que nunca haya saldos ni libros inconsistentes}{que el motor rechace cualquier operación que rompa un invariante y que se prueben con datos aleatorios}}
\begin{gherkin}
\kw{Característica:} Invariantes del motor
\nv{1}\kw{Escenario:} La distribución suma el abono
\nv{2}\kw{Dado} cualquier abono aplicado
\nv{2}\kw{Cuando} se suma su distribución por concepto
\nv{2}\kw{Entonces} el total es igual al valor del abono
\nv{1}\kw{Escenario:} Crédito cerrado
\nv{2}\kw{Dado} un crédito en estado "Cancelado" o "Cerrado por evento"
\nv{2}\kw{Cuando} llega un abono
\nv{2}\kw{Entonces} el motor lo rechaza y no modifica saldos
\nv{1}\kw{Escenario:} Pruebas con datos aleatorios
\nv{2}\kw{Dado} 10.000 secuencias aleatorias de abonos válidos
\nv{2}\kw{Cuando} se aplican
\nv{2}\kw{Entonces} ningún saldo queda negativo y todos los asientos cuadran
\end{gherkin}
\nota{en el Sprint~1 se cubre la parte del cálculo: el saldo nunca es negativo y el cronograma
cierra en cero para cualquier monto, tasa y plazo válidos. Los escenarios de abonos y asientos se
completan en los sprints 2 y 3.}

% ═════════════════════════════════════════════════════════════════════════════
\section{Historias nuevas de los sprints 2 y 3}

\hu{HU-07}{Contrato de datos con el CRM}{\must}{E6 · Integración \textit{(nueva, Sprint 2)}}{%
\como{área de cartera}{que el CRM y el motor acuerden qué datos se mandan, qué se recibe y quién es dueño de cada dato}{que los abonos lleguen completos y no se dupliquen gestiones}}
\begin{gherkin}
\kw{Característica:} Contrato de integración CRM–motor
\nv{1}\kw{Escenario:} Contrato acordado
\nv{2}\kw{Dado} el CRM que registra los abonos de los clientes
\nv{2}\kw{Cuando} se firma el contrato de integración
\nv{2}\kw{Entonces} quedan definidos los campos del abono, el identificador único,
\nv{2}el formato de respuesta, los avisos que recibe el CRM y los códigos de error
\nv{2}\kw{Y} queda definido que el motor es dueño de saldos, estados y asientos
\nv{2}\kw{Y} que el CRM es dueño de clientes, contactos y gestiones de cobro
\end{gherkin}

\hu{HU-19}{Máquina de estados del crédito}{\must}{E3 · Eventos \textit{(nueva, Sprint 2)}}{%
\como{analista de cartera}{que el crédito solo cambie de estado por las transiciones permitidas}{que el estado siempre refleje la realidad del crédito}}
\begin{gherkin}
\kw{Característica:} Estados del crédito
\nv{1}\kw{Esquema del escenario:} Transición permitida
\nv{2}\kw{Dado} un crédito en estado "<desde>"
\nv{2}\kw{Cuando} ocurre "<evento>"
\nv{2}\kw{Entonces} el crédito queda en "<hacia>"
\nv{2}\kw{Ejemplos:}
\nv{3}| Vigente al día | pago inferior \ \ \ \ \ \ | En mora \ \ \ \ \ \ \ \ \ \ \ |
\nv{3}| En mora \ \ \ \ \ \ \ | cubre todo lo vencido | Vigente al día \ \ \ \ |
\nv{3}| Vigente al día | prepago total \ \ \ \ \ \ | Cancelado \ \ \ \ \ \ \ \ \ |
\nv{3}| En mora \ \ \ \ \ \ \ | dación o retoma \ \ \ \ | Cerrado por evento |
\nv{1}\kw{Escenario:} Transición no permitida
\nv{2}\kw{Dado} un crédito "Cancelado"
\nv{2}\kw{Cuando} se intenta pasarlo a "En mora"
\nv{2}\kw{Entonces} el motor lo rechaza y registra el intento
\end{gherkin}

\hu{HT-03}{Historial inmutable de movimientos y recálculo}{\must}{E9 · Plataforma \textit{(técnica, Sprint 2)}}{%
\paraht{reconstruir el crédito a cualquier fecha y aceptar pagos con fecha atrasada}{guardar cada abono, reverso o cambio de tasa como un movimiento que no se borra y recalcular recorriéndolos en orden}}
\begin{gherkin}
\kw{Característica:} Historial de movimientos
\nv{1}\kw{Escenario:} Pago registrado con fecha atrasada
\nv{2}\kw{Dado} un crédito con abonos aplicados hasta la cuota 45
\nv{2}\kw{Cuando} llega un abono con fecha de la cuota 43 que no se había registrado
\nv{2}\kw{Entonces} el motor recalcula el crédito desde la cuota 43 en orden cronológico
\nv{2}\kw{Y} genera los asientos de ajuste necesarios
\nv{1}\kw{Escenario:} Consulta a una fecha
\nv{2}\kw{Dado} el mismo crédito
\nv{2}\kw{Cuando} se consulta su estado al cierre de la cuota 40
\nv{2}\kw{Entonces} se obtienen los saldos tal como estaban en esa fecha
\end{gherkin}

\hu{HT-07}{Avisos al CRM sin pérdida (outbox)}{\must}{E6 · Integración \textit{(técnica, Sprint 3)}}{%
\paraht{que el CRM no pierda ningún cambio de estado aunque esté caído}{guardar cada aviso en la misma transacción del abono y enviarlo con reintentos}}
\begin{gherkin}
\kw{Característica:} Envío confiable de avisos
\nv{1}\kw{Escenario:} El CRM no responde
\nv{2}\kw{Dado} un abono que deja el crédito "En mora"
\nv{2}\kw{Y} el CRM caído
\nv{2}\kw{Cuando} se aplica el abono
\nv{2}\kw{Entonces} el aviso queda guardado en la misma transacción del abono
\nv{2}\kw{Y} se reintenta hasta que el CRM confirme la recepción
\nv{2}\kw{Y} el CRM no recibe el mismo aviso dos veces como si fuera nuevo
\end{gherkin}

\hu{HT-11}{Ambiente de pruebas desplegado}{\must}{E9 · Plataforma \textit{(técnica, nueva, Sprint 3)}}{%
\paraht{probar la integración con el CRM antes de producción}{un ambiente de pruebas con su base de datos que se actualiza solo cuando se une un cambio a la rama principal}}
\begin{gherkin}
\kw{Característica:} Ambiente de pruebas
\nv{1}\kw{Escenario:} Despliegue automático
\nv{2}\kw{Dado} un cambio unido a la rama principal con el pipeline en verde
\nv{2}\kw{Cuando} termina el pipeline
\nv{2}\kw{Entonces} la nueva versión queda desplegada en el ambiente de pruebas
\nv{2}\kw{Y} el CRM de pruebas puede llamar a la API con sus credenciales
\nv{1}\kw{Escenario:} Datos separados
\nv{2}\kw{Dado} el ambiente de pruebas
\nv{2}\kw{Cuando} se revisa su configuración
\nv{2}\kw{Entonces} usa su propia base de datos y ninguna credencial de producción
\end{gherkin}

% ═════════════════════════════════════════════════════════════════════════════
\section{Entrega (handoff)}
\begin{itemize}
  \item \textbf{¿Listo para implementar?} Sí el Sprint~1: base técnica (HT-08, HT-01, HT-09, HT-10) y núcleo del caso~1. Los casos 8 y 9 esperan P1, P3 y P4.
  \item \textbf{¿Necesita revisión de arquitectura?} Sí, confirmar el ADR del monolito modular (HT-01) y, antes del Sprint~2, el contrato con el CRM (HU-07).
  \item \textbf{¿Necesita aclaración de producto?} Sí: preguntas P1 a P10 con Diego González y Heider Galvis.
  \item \textbf{Siguiente paso:} el lunes del Sprint~1, sesión de refinamiento de 2 horas con negocio,
        repositorio y ADR; desde el martes, núcleo de cálculo con pruebas primero (TDD) contra la hoja \textit{Validaciones}.
\end{itemize}

\end{document}
